\section{总结与延伸}

\subsection{全讲框架}

前面从评估基础走到生产治理，本节将任务、环境、测量、项目与维护压缩成一张总表。读表时应把每行视为必需 artifact，而不是可选文档；任何一层缺失，都可能让高分失去可复现性或真实含义。

\begin{longtable}{p{2.7cm}p{4.3cm}p{4.1cm}}
\toprule
\textbf{层次} & \textbf{核心问题} & \textbf{交付物} \\
\midrule
\endfirsthead
\toprule
\textbf{层次} & \textbf{核心问题} & \textbf{交付物} \\
\midrule
\endhead
任务 & 要测什么，是否有 oracle & task spec、scope、constraints \\
环境 & Agent 在哪里行动 & versioned sandbox、tool contract \\
评估 & 分数是否代表真实 outcome & oracle/judge、validity audit \\
数据 & task 如何生成和维护 & pipeline、provenance、split \\
可靠性 & 重复运行是否稳定 & variance、failure taxonomy \\
项目 & benchmark 如何作为服务运行 & Green Agent、protocol、report \\
治理 & 是否安全、可复现、可审计 & permission、digest、trace \\
\bottomrule
\end{longtable}

\subsection{十二条核心结论}

\begin{importantbox}{从 Eval 到 Green Agent}
\begin{enumerate}[leftmargin=*]
\item Evaluation 是系统性、可重复的决策接口。
\item Agent eval 的单位是 trajectory 与 environment outcome。
\item Close/open 与 verifiable/non-verifiable 是不同维度。
\item 能用确定性 oracle 时，不应只依赖 LLM judge。
\item Human 和 LLM judge 都需要校准、一致性与偏差报告。
\item Static core set 与 dynamic rotating set 应双轨运行。
\item Capability、application、general-set eval 各有诊断职责。
\item Good eval 的核心是 outcome validity。
\item Benchmark 必须公开 data pipeline、限制与污染风险。
\item Green Agent 是 evaluator/orchestrator，White Agent 是 participant。
\item Existing integration 也必须做质量分析和 baseline parity。
\item 可复现镜像、协议、trace 与 verifier evidence 是项目的一部分。
\end{enumerate}
\end{importantbox}

\subsection{拓展阅读}

\begin{itemize}
\item Survey on Evaluation of LLM-based Agents：能力、应用与通用评估 taxonomy。
\item Outcome validity 相关工作：结果是否与真实任务意图一致。
\item CyberGym、tau-bench/tau2-bench、GDPval、CRMArena、LegalAgentBench。
\item SWE-bench Verified：benchmark 任务审查与可靠性提升。
\item AgentBeats、A2A 与 Google ADK：多 Agent 发现、通信与运行工具。
\end{itemize}

\subsection{开放问题}

当环境无法完全复现时，怎样报告 uncertainty？动态 benchmark 如何保持跨版本可比？Green Agent 如何抵抗 participant prompt injection？如何把安全违规设为硬门槛，同时保留多目标 Pareto 分析？这些问题说明 Agent evaluation 不是 leaderboard 的附属工作，而是未来 Agent 工程和科学方法的共同基础。

\end{document}
